\documentclass[11pt,letterpaper]{article}

% Paquetes necesarios
\usepackage[utf8]{inputenc}
\usepackage[spanish]{babel}
\usepackage[margin=2cm]{geometry}
\usepackage{fancyhdr}
\usepackage{graphicx}
\usepackage{listings}
\usepackage{xcolor}
\usepackage{hyperref}
\usepackage{tabularx}
\usepackage{array}

% Configuración del header y footer
\pagestyle{fancy}
\fancyhf{}
\lhead{NLPY-1: Fundamentos de Python}
\rhead{Semana 7: Expresiones Regulares}
\rfoot{Página \thepage}

% Configuración de listings para Python
\lstset{
    language=Python,
    basicstyle=\ttfamily\small,
    keywordstyle=\color{blue}\bfseries,
    stringstyle=\color{red},
    commentstyle=\color{gray}\itshape,
    numbers=left,
    numberstyle=\tiny\color{gray},
    stepnumber=1,
    numbersep=8pt,
    showstringspaces=false,
    breaklines=true,
    frame=single,
    backgroundcolor=\color{gray!10},
    tabsize=4,
    captionpos=b
}

% Configuración de hyperref
\hypersetup{
    colorlinks=true,
    linkcolor=blue,
    urlcolor=blue,
    citecolor=blue
}

\begin{document}

% Título
\begin{center}
    {\LARGE \textbf{Tarea 7: Validador y Extractor de Datos}}\\[0.3cm]
    {\Large Semana 7: Expresiones Regulares (RegEx)}\\[0.2cm]
    {\large NLPY-1: Fundamentos de Python}\\[0.2cm]
\end{center}

\section{Introducción}

Las \textbf{expresiones regulares} (regular expressions o regex) son una de las herramientas más poderosas en la programación moderna. Permiten buscar, validar y extraer patrones de texto de manera eficiente y elegante. Son ampliamente utilizadas en validación de formularios, procesamiento de logs, web scraping, análisis de datos y mucho más.

En esta semana aprenderás a dominar las expresiones regulares en Python utilizando el módulo \texttt{re}. Desarrollarás un sistema completo de validación y extracción de datos que podrá procesar información de diferentes formatos y fuentes.

\subsection{¿Por qué son importantes las Expresiones Regulares?}

\begin{itemize}
    \item \textbf{Validación de datos}: Verificar que emails, teléfonos, cédulas, etc. tengan el formato correcto
    \item \textbf{Extracción de información}: Obtener datos específicos de textos grandes o documentos
    \item \textbf{Búsqueda avanzada}: Encontrar patrones complejos que serían difíciles con métodos tradicionales
    \item \textbf{Transformación de texto}: Reemplazar o reformatear información de manera masiva
    \item \textbf{Parsing de logs}: Analizar archivos de registro de sistemas y aplicaciones
\end{itemize}

\subsection{Objetivos de Aprendizaje}

Al completar esta tarea, el estudiante será capaz de:

\begin{itemize}
    \item Comprender la sintaxis y estructura de las expresiones regulares
    \item Utilizar metacaracteres, cuantificadores y clases de caracteres
    \item Aplicar las funciones del módulo \texttt{re}: \texttt{match()}, \texttt{search()}, \texttt{findall()}, \texttt{sub()}
    \item Crear patrones para validar formatos costarricenses (cédulas, teléfonos, placas)
    \item Usar grupos de captura y grupos con nombre
    \item Implementar flags para modificar el comportamiento de búsqueda
    \item Extraer información estructurada de texto no estructurado
    \item Aplicar expresiones regulares para limpieza y sanitización de datos
\end{itemize}

\section{Descripción del Proyecto}

Desarrollarás un \textbf{Sistema de Validación y Extracción de Datos} que procese diferentes tipos de información usando expresiones regulares. El sistema debe validar formatos costarricenses, extraer datos de textos, y generar reportes estructurados.

\subsection{Funcionalidades Requeridas}

\subsubsection{1. Validadores (20 puntos - Funcionalidad)}

Implementar validadores para formatos costarricenses usando regex:

\textbf{A) Validador de Cédulas (5 puntos)}
\begin{itemize}
    \item Cédula física: 1-2345-6789 (formato: D-DDDD-DDDD)
    \item Cédula jurídica: 3-101-123456 (formato: D-DDD-DDDDDD)
    \item DIMEX: 123456789012 (12 dígitos)
    \item NITE: 1234567890 (10 dígitos)
    \item Validar formato y estructura correcta
\end{itemize}

\textbf{B) Validador de Teléfonos (4 puntos)}
\begin{itemize}
    \item Formato nacional: 8888-8888, 2222-2222
    \item Con código de país: +506 8888-8888, +506-8888-8888
    \item Sin guiones: 88888888
    \item Validar que inicie con 2, 4, 5, 6, 7 u 8
\end{itemize}

\textbf{C) Validador de Emails (3 puntos)}
\begin{itemize}
    \item Formato estándar: usuario@dominio.com
    \item Permitir puntos, guiones y guiones bajos en usuario
    \item Validar dominio y extensión
    \item Soportar subdominios: usuario@mail.empresa.co.cr
\end{itemize}

\textbf{D) Validador de Placas de Vehículos (4 puntos)}
\begin{itemize}
    \item Placas particulares: ABC-123, ABC-1234
    \item Placas comerciales: C-12345
    \item Placas públicas: PA-1234, TP-12345
    \item Placas de motos: A-12345
    \item Case-insensitive
\end{itemize}

\textbf{E) Validador de URLs (4 puntos)}
\begin{itemize}
    \item HTTP/HTTPS: https://www.ejemplo.com
    \item Con rutas: https://ejemplo.com/ruta/archivo.html
    \item Con parámetros: https://ejemplo.com?param=valor\&otro=123
    \item Con puertos: https://ejemplo.com:8080
\end{itemize}

\subsubsection{2. Extractores (20 puntos - Funcionalidad)}

Implementar extractores que obtengan información de textos:

\textbf{A) Extractor de Información Personal (8 puntos)}
\begin{itemize}
    \item Extraer todos los emails de un texto
    \item Extraer todos los teléfonos
    \item Extraer todas las cédulas
    \item Extraer fechas (formatos: DD/MM/YYYY, DD-MM-YYYY, YYYY-MM-DD)
    \item Retornar listas con los datos encontrados
\end{itemize}

\textbf{B) Extractor de Precios y Montos (4 puntos)}
\begin{itemize}
    \item Identificar montos en colones: ₡12,345.50, ₡1.500
    \item Identificar montos en dólares: \$1,234.56
    \item Extraer números con separadores de miles
    \item Capturar valor numérico y moneda
\end{itemize}

\textbf{C) Procesador de Logs (8 puntos)}
\begin{itemize}
    \item Procesar líneas de log de formato estándar
    \item Extraer: timestamp, nivel (INFO/ERROR/WARNING), mensaje
    \item Ejemplo: \texttt{[2024-10-24 14:30:15] ERROR: Connection failed}
    \item Generar diccionario estructurado por cada entrada
    \item Contar ocurrencias por nivel de severidad
\end{itemize}

\subsubsection{3. Requisitos Técnicos}

\textbf{Uso del Módulo re:}
\begin{itemize}
    \item Usar \texttt{re.compile()} para patrones reutilizables
    \item Aplicar \texttt{re.match()} para validación de inicio
    \item Usar \texttt{re.search()} para encontrar patrones en texto
    \item Aplicar \texttt{re.findall()} para extraer múltiples coincidencias
    \item Usar \texttt{re.sub()} para reemplazos y limpieza
\end{itemize}

\textbf{Características Avanzadas:}
\begin{itemize}
    \item Implementar grupos de captura con paréntesis
    \item Usar grupos con nombre: \texttt{(?P<nombre>...)}
    \item Aplicar flags: \texttt{re.IGNORECASE}, \texttt{re.MULTILINE}, \texttt{re.DOTALL}
    \item Implementar lookahead y lookbehind (opcional)
\end{itemize}

\textbf{Organización del Código:}
\begin{itemize}
    \item Función validadora por cada tipo de dato
    \item Función extractora que retorne listas o diccionarios
    \item Diccionario global con patrones regex compilados
    \item Menú interactivo para probar validadores
    \item Función de reporte que muestre estadísticas
\end{itemize}

\section{Conceptos de Expresiones Regulares}

\subsection{Metacaracteres Básicos}

\begin{itemize}
    \item \texttt{.} - Cualquier carácter (excepto salto de línea)
    \item \texttt{\^{}} - Inicio de cadena o línea
    \item \texttt{\$} - Fin de cadena o línea
    \item \texttt{*} - 0 o más repeticiones
    \item \texttt{+} - 1 o más repeticiones
    \item \texttt{?} - 0 o 1 repetición (opcional)
    \item \texttt{\{n\}} - Exactamente n repeticiones
    \item \texttt{\{n,m\}} - Entre n y m repeticiones
    \item \texttt{|} - OR lógico
    \item \texttt{[]} - Clase de caracteres
    \item \texttt{()} - Grupo de captura
\end{itemize}

\subsection{Clases de Caracteres}

\begin{itemize}
    \item \texttt{[abc]} - a, b, o c
    \item \texttt{[a-z]} - Cualquier letra minúscula
    \item \texttt{[A-Z]} - Cualquier letra mayúscula
    \item \texttt{[0-9]} - Cualquier dígito
    \item \texttt{[\^{}abc]} - Cualquier cosa EXCEPTO a, b, c
    \item \texttt{\textbackslash d} - Dígito [0-9]
    \item \texttt{\textbackslash D} - No dígito
    \item \texttt{\textbackslash w} - Alfanumérico [a-zA-Z0-9\_]
    \item \texttt{\textbackslash W} - No alfanumérico
    \item \texttt{\textbackslash s} - Espacio en blanco
    \item \texttt{\textbackslash S} - No espacio
\end{itemize}

\subsection{Ejemplos de Patrones}

\begin{lstlisting}[caption={Ejemplos de patrones regex comunes}]
# Cedula fisica costarricense: 1-2345-6789
patron_cedula = r'^\d-\d{4}-\d{4}$'

# Telefono: 8888-8888 o 88888888
patron_telefono = r'^[2-8]\d{3}-?\d{4}$'

# Email basico
patron_email = r'^[a-zA-Z0-9._-]+@[a-zA-Z0-9.-]+\.[a-zA-Z]{2,}$'

# Placa particular: ABC-123 o ABC-1234
patron_placa = r'^[A-Z]{3}-\d{3,4}$'

# Fecha DD/MM/YYYY
patron_fecha = r'\d{2}/\d{2}/\d{4}'

# Monto en colones: ₡12,345.50
patron_colones = r'₡[\d,]+\.?\d{0,2}'

# Log entry
patron_log = r'\[(\d{4}-\d{2}-\d{2} \d{2}:\d{2}:\d{2})\] (\w+): (.+)'
\end{lstlisting}

\section{Plantilla de Código}

\subsection{Estructura Base del Sistema}

\begin{lstlisting}[caption={Plantilla del validador y extractor}]
import re
from typing import List, Dict, Optional

# ==================================================================
# PATRONES REGEX COMPILADOS
# ==================================================================

PATRONES = {
    # Cedulas
    'cedula_fisica': re.compile(r''),  # TODO: patron para 1-2345-6789
    'cedula_juridica': re.compile(r''),  # TODO: patron para 3-101-123456
    'dimex': re.compile(r''),  # TODO: 12 digitos
    'nite': re.compile(r''),  # TODO: 10 digitos
    
    # Telefonos
    'telefono': re.compile(r''),  # TODO: patrones variados de telefono
    
    # Email
    'email': re.compile(r''),  # TODO: patron de email
    
    # Placas
    'placa_particular': re.compile(r'', re.IGNORECASE),  # TODO: ABC-123
    'placa_comercial': re.compile(r'', re.IGNORECASE),  # TODO: C-12345
    'placa_moto': re.compile(r'', re.IGNORECASE),  # TODO: A-12345
    
    # URLs
    'url': re.compile(r''),  # TODO: http(s)://...
    
    # Fechas
    'fecha_slash': re.compile(r''),  # TODO: DD/MM/YYYY
    'fecha_guion': re.compile(r''),  # TODO: DD-MM-YYYY
    'fecha_iso': re.compile(r''),  # TODO: YYYY-MM-DD
    
    # Montos
    'colones': re.compile(r''),  # TODO: ₡12,345.50
    'dolares': re.compile(r''),  # TODO: $1,234.56
    
    # Logs
    'log_entry': re.compile(r'')  # TODO: [timestamp] NIVEL: mensaje
}

# ==================================================================
# VALIDADORES
# ==================================================================

def validar_cedula(cedula: str) -> Dict[str, any]:
    """
    Valida una cedula costarricense (fisica, juridica, DIMEX, NITE).
    
    Returns:
        dict: {'valida': bool, 'tipo': str, 'mensaje': str}
    """
    # TODO: Probar contra cada patron de cedula
    # TODO: Retornar tipo identificado
    pass

def validar_telefono(telefono: str) -> Dict[str, any]:
    """
    Valida un numero de telefono costarricense.
    
    Returns:
        dict: {'valido': bool, 'formato': str, 'mensaje': str}
    """
    # TODO: Validar patron de telefono
    # TODO: Verificar que inicie con digito valido
    pass

def validar_email(email: str) -> Dict[str, any]:
    """
    Valida un email.
    
    Returns:
        dict: {'valido': bool, 'usuario': str, 'dominio': str}
    """
    # TODO: Validar patron de email
    # TODO: Extraer usuario y dominio usando grupos
    pass

def validar_placa(placa: str) -> Dict[str, any]:
    """
    Valida una placa vehicular costarricense.
    
    Returns:
        dict: {'valida': bool, 'tipo': str, 'mensaje': str}
    """
    # TODO: Validar contra patrones de placas
    # TODO: Identificar tipo de placa
    pass

def validar_url(url: str) -> Dict[str, any]:
    """
    Valida una URL.
    
    Returns:
        dict: {'valida': bool, 'protocolo': str, 'dominio': str}
    """
    # TODO: Validar patron de URL
    # TODO: Extraer componentes con grupos
    pass

# ==================================================================
# EXTRACTORES
# ==================================================================

def extraer_emails(texto: str) -> List[str]:
    """Extrae todos los emails de un texto."""
    # TODO: Usar findall para encontrar todos los emails
    pass

def extraer_telefonos(texto: str) -> List[str]:
    """Extrae todos los telefonos de un texto."""
    # TODO: Usar findall para encontrar todos los telefonos
    pass

def extraer_cedulas(texto: str) -> List[Dict[str, str]]:
    """Extrae todas las cedulas de un texto."""
    # TODO: Buscar todos los tipos de cedula
    # TODO: Retornar lista con tipo y valor
    pass

def extraer_fechas(texto: str) -> List[str]:
    """Extrae todas las fechas de un texto."""
    # TODO: Buscar todos los formatos de fecha
    pass

def extraer_montos(texto: str) -> List[Dict[str, any]]:
    """
    Extrae montos con moneda de un texto.
    
    Returns:
        list: [{'monto': str, 'moneda': str, 'valor': float}, ...]
    """
    # TODO: Extraer montos en colones y dolares
    # TODO: Convertir a valores numericos
    pass

def extraer_datos_personales(texto: str) -> Dict[str, List]:
    """
    Extrae toda la informacion personal de un texto.
    
    Returns:
        dict: {
            'emails': [...],
            'telefonos': [...],
            'cedulas': [...],
            'fechas': [...]
        }
    """
    # TODO: Usar todas las funciones extractoras
    # TODO: Retornar diccionario organizado
    pass

# ==================================================================
# PROCESADOR DE LOGS
# ==================================================================

def procesar_log(texto_log: str) -> Dict[str, any]:
    """
    Procesa un archivo de log y extrae informacion estructurada.
    
    Returns:
        dict: {
            'entradas': [...],
            'estadisticas': {'INFO': n, 'ERROR': n, ...}
        }
    """
    # TODO: Extraer cada linea de log
    # TODO: Parsear timestamp, nivel, mensaje
    # TODO: Contar por nivel de severidad
    pass

# ==================================================================
# UTILIDADES
# ==================================================================

def limpiar_telefono(telefono: str) -> str:
    """Limpia un telefono dejando solo digitos."""
    # TODO: Usar re.sub para eliminar no-digitos
    pass

def formatear_cedula(cedula: str) -> str:
    """Formatea una cedula al formato estandar."""
    # TODO: Extraer digitos y reformatear
    pass

def enmascarar_email(email: str) -> str:
    """Enmascara un email: usuario@dominio.com -> u***o@dominio.com"""
    # TODO: Usar grupos para capturar y reemplazar
    pass

# ==================================================================
# MENU INTERACTIVO
# ==================================================================

def menu_validadores():
    """Menu para probar los validadores."""
    while True:
        print("\n" + "="*50)
        print("VALIDADORES")
        print("="*50)
        print("1. Validar cedula")
        print("2. Validar telefono")
        print("3. Validar email")
        print("4. Validar placa")
        print("5. Validar URL")
        print("6. Volver")
        
        # TODO: Implementar logica del menu
        pass

def menu_extractores():
    """Menu para probar los extractores."""
    # TODO: Implementar menu de extractores
    pass

def menu_principal():
    """Menu principal del sistema."""
    while True:
        print("\n" + "="*60)
        print("SISTEMA DE VALIDACION Y EXTRACCION DE DATOS")
        print("="*60)
        print("1. Validadores")
        print("2. Extractores de texto")
        print("3. Procesador de logs")
        print("4. Utilidades")
        print("5. Ejecutar pruebas")
        print("6. Salir")
        
        # TODO: Implementar navegacion del menu
        pass

# ==================================================================
# PRUEBAS
# ==================================================================

def ejecutar_pruebas():
    """Ejecuta casos de prueba para validar el sistema."""
    print("\n" + "="*60)
    print("EJECUTANDO PRUEBAS...")
    print("="*60)
    
    # Pruebas de cedulas
    cedulas_prueba = [
        "1-2345-6789",      # Fisica valida
        "3-101-123456",     # Juridica valida
        "123456789012",     # DIMEX valido
        "1-234-5678",       # Fisica invalida
    ]
    
    # TODO: Ejecutar validaciones y mostrar resultados
    # TODO: Probar extractores con textos de ejemplo
    # TODO: Mostrar estadisticas de exito/fallo
    pass

# ==================================================================
# MAIN
# ==================================================================

if __name__ == "__main__":
    print("Sistema de Validacion y Extraccion de Datos")
    print("NLPY-1 - Semana 7: Expresiones Regulares")
    menu_principal()
\end{lstlisting}

\subsection{Ejemplo: Implementación de un Validador}

Como referencia, aquí hay un ejemplo de cómo implementar el validador de email:

\begin{lstlisting}[caption={Ejemplo: Validador de email completo}]
def validar_email(email: str) -> Dict[str, any]:
    """
    Valida un email y extrae sus componentes.
    
    Args:
        email: String con el email a validar
        
    Returns:
        dict: {
            'valido': bool,
            'usuario': str,
            'dominio': str,
            'mensaje': str
        }
    """
    # Patron con grupos de captura
    patron = re.compile(
        r'^(?P<usuario>[a-zA-Z0-9._-]+)@(?P<dominio>[a-zA-Z0-9.-]+\.[a-zA-Z]{2,})$'
    )
    
    match = patron.match(email.strip())
    
    if match:
        return {
            'valido': True,
            'usuario': match.group('usuario'),
            'dominio': match.group('dominio'),
            'mensaje': 'Email valido'
        }
    else:
        return {
            'valido': False,
            'usuario': None,
            'dominio': None,
            'mensaje': 'Formato de email invalido'
        }

# Ejemplos de uso:
print(validar_email("juan.perez@tec.ac.cr"))
# {'valido': True, 'usuario': 'juan.perez', 'dominio': 'tec.ac.cr', ...}

print(validar_email("correo_invalido@"))
# {'valido': False, 'usuario': None, 'dominio': None, ...}
\end{lstlisting}

\section{Ejemplos de Ejecución}

\subsection{Validación de Cédula}

\begin{lstlisting}[caption={Ejemplo de validación de cédulas}]
>>> validar_cedula("1-2345-6789")
{
    'valida': True,
    'tipo': 'Fisica',
    'mensaje': 'Cedula fisica valida'
}

>>> validar_cedula("3-101-654321")
{
    'valida': True,
    'tipo': 'Juridica',
    'mensaje': 'Cedula juridica valida'
}

>>> validar_cedula("1-234-567")
{
    'valida': False,
    'tipo': None,
    'mensaje': 'Formato de cedula invalido'
}
\end{lstlisting}

\subsection{Extracción de Datos}

\begin{lstlisting}[caption={Ejemplo de extracción de datos personales}]
texto = """
Contactar a Juan Perez al telefono 8888-8888 o al email 
juan.perez@tec.ac.cr. Cedula: 1-2345-6789.
Tambien puede llamar a Maria al 2222-2222 o escribir a 
maria.rodriguez@email.com. Fecha de nacimiento: 15/03/1990.
"""

>>> extraer_datos_personales(texto)
{
    'emails': ['juan.perez@tec.ac.cr', 'maria.rodriguez@email.com'],
    'telefonos': ['8888-8888', '2222-2222'],
    'cedulas': [{'tipo': 'fisica', 'valor': '1-2345-6789'}],
    'fechas': ['15/03/1990']
}
\end{lstlisting}

\subsection{Procesamiento de Logs}

\begin{lstlisting}[caption={Ejemplo de procesamiento de logs}]
log_texto = """
[2024-10-24 14:30:15] INFO: Sistema iniciado correctamente
[2024-10-24 14:30:16] INFO: Conectando a base de datos
[2024-10-24 14:30:20] ERROR: Conexion fallida, reintentando
[2024-10-24 14:30:25] WARNING: Tiempo de respuesta elevado
[2024-10-24 14:30:30] INFO: Conexion establecida
"""

>>> procesar_log(log_texto)
{
    'entradas': [
        {
            'timestamp': '2024-10-24 14:30:15',
            'nivel': 'INFO',
            'mensaje': 'Sistema iniciado correctamente'
        },
        ...
    ],
    'estadisticas': {
        'INFO': 3,
        'ERROR': 1,
        'WARNING': 1
    },
    'total': 5
}
\end{lstlisting}

\section{Defensa del Código (60\% de la Nota)}

\subsection{Aspectos Técnicos a Dominar}

\textbf{1. Sintaxis de Expresiones Regulares (15 puntos)}
\begin{itemize}
    \item Explicar metacaracteres: \texttt{.}, \texttt{*}, \texttt{+}, \texttt{?}, \texttt{\^{}}, \texttt{\$}
    \item ¿Cuál es la diferencia entre \texttt{*} y \texttt{+}? ¿Y entre \texttt{?} y \texttt{\{0,1\}}?
    \item ¿Qué hacen los cuantificadores \texttt{\{n\}}, \texttt{\{n,\}}, \texttt{\{n,m\}}?
    \item Explicar clases de caracteres: \texttt{[abc]}, \texttt{[a-z]}, \texttt{[\^{}abc]}
    \item Diferencia entre \texttt{\textbackslash d} y \texttt{[0-9]}
\end{itemize}

\textbf{2. Módulo re de Python (15 puntos)}
\begin{itemize}
    \item Diferencia entre \texttt{re.match()} y \texttt{re.search()}
    \item ¿Cuándo usar \texttt{re.findall()} vs \texttt{re.finditer()}?
    \item ¿Para qué sirve \texttt{re.compile()}? ¿Ventajas de compilar patrones?
    \item Explicar el uso de \texttt{re.sub()} para reemplazos
    \item ¿Qué retorna cada función? (objeto Match, lista, string, etc.)
\end{itemize}

\textbf{3. Grupos de Captura (10 puntos)}
\begin{itemize}
    \item ¿Qué es un grupo de captura? ¿Cómo se define?
    \item Diferencia entre grupos normales y grupos con nombre
    \item ¿Cómo acceder a grupos capturados? (\texttt{.group()}, \texttt{.groups()})
    \item ¿Qué son los grupos no capturadores \texttt{(?:...)}?
    \item Explicar backreferences
\end{itemize}

\textbf{4. Flags y Modificadores (10 puntos)}
\begin{itemize}
    \item ¿Para qué sirve \texttt{re.IGNORECASE}?
    \item Explicar \texttt{re.MULTILINE} y su efecto en \texttt{\^{}} y \texttt{\$}
    \item ¿Qué hace \texttt{re.DOTALL}?
    \item ¿Cómo combinar múltiples flags?
    \item ¿Qué es \texttt{re.VERBOSE} y cuándo usarlo?
\end{itemize}

\textbf{5. Aplicaciones Prácticas (10 puntos)}
\begin{itemize}
    \item ¿Por qué validar con regex en lugar de métodos string?
    \item Limitaciones de las expresiones regulares
    \item Casos donde NO se deben usar regex
    \item Rendimiento: ¿cuándo compilar patrones?
    \item Mejores prácticas y patrones comunes
\end{itemize}

\subsection{Preguntas de Defensa Típicas}

\begin{itemize}
    \item "Explica línea por línea tu patrón regex para validar cédulas físicas"
    \item "¿Por qué usaste \texttt{re.match()} en lugar de \texttt{re.search()}?"
    \item "Muestra cómo extraerías el código de área de un teléfono usando grupos"
    \item "¿Qué pasa si no usas \texttt{r''} (raw string) en tus patrones?"
    \item "Modifica tu patrón de email para que también acepte signos de suma en el usuario"
    \item "¿Cómo harías para que tu validador de placas sea case-insensitive?"
    \item "Explica cómo funciona tu extractor de logs paso a paso"
\end{itemize}

\section{Rúbrica de Evaluación}

\subsection{Funcionalidad del Código (40 puntos)}

\textbf{Validadores (20 puntos):}
\begin{itemize}
    \item 5 puntos: Validador de cédulas (todos los tipos)
    \item 4 puntos: Validador de teléfonos (formatos variados)
    \item 3 puntos: Validador de emails
    \item 4 puntos: Validador de placas (todos los tipos)
    \item 4 puntos: Validador de URLs
\end{itemize}

\textbf{Extractores (20 puntos):}
\begin{itemize}
    \item 8 puntos: Extractor de datos personales completo
    \item 4 puntos: Extractor de montos y moneda
    \item 8 puntos: Procesador de logs con estadísticas
\end{itemize}

\subsection{Defensa del Código (60 puntos)}

\textbf{Sintaxis de Regex (15 puntos):}
\begin{itemize}
    \item 6 puntos: Explica metacaracteres y cuantificadores
    \item 5 puntos: Comprende clases de caracteres
    \item 4 puntos: Entiende anclas y límites de palabra
\end{itemize}

\textbf{Módulo re (15 puntos):}
\begin{itemize}
    \item 5 puntos: Diferencia entre funciones principales
    \item 5 puntos: Uso correcto de re.compile()
    \item 5 puntos: Comprende objetos Match y sus métodos
\end{itemize}

\textbf{Grupos y Captura (10 puntos):}
\begin{itemize}
    \item 5 puntos: Implementa grupos de captura correctamente
    \item 3 puntos: Usa grupos con nombre apropiadamente
    \item 2 puntos: Accede a grupos capturados
\end{itemize}

\textbf{Flags y Técnicas Avanzadas (10 puntos):}
\begin{itemize}
    \item 4 puntos: Aplica flags correctamente
    \item 3 puntos: Usa raw strings
    \item 3 puntos: Implementa técnicas avanzadas
\end{itemize}

\textbf{Calidad y Organización (10 puntos):}
\begin{itemize}
    \item 3 puntos: Patrones bien documentados
    \item 3 puntos: Código organizado y limpio
    \item 2 puntos: Manejo de errores
    \item 2 puntos: Pruebas y ejemplos
\end{itemize}

\subsection{Penalizaciones}

\begin{itemize}
    \item \textbf{-5 puntos}: No usar raw strings (\texttt{r''}) en patrones regex
    \item \textbf{-5 puntos}: No usar \texttt{re.compile()} para patrones reutilizables
    \item \textbf{-3 puntos}: Validadores que no manejan casos límite
    \item \textbf{-3 puntos}: No usar grupos de captura apropiadamente
    \item \textbf{-2 puntos}: Patrones regex sin comentarios explicativos
    \item \textbf{-2 puntos}: No validar entrada antes de procesar
    \item \textbf{-5 puntos}: Código copiado sin comprensión
\end{itemize}

\section{Desafíos Opcionales (Puntos Extra)}

\subsection{Nivel Intermedio (+10 puntos cada uno)}

\textbf{1. Validador de Contraseñas Seguras}
\begin{itemize}
    \item Mínimo 8 caracteres
    \item Al menos una mayúscula, minúscula, número y símbolo
    \item Usar lookahead assertions: \texttt{(?=.*[A-Z])}
\end{itemize}

\textbf{2. Extractor de Hashtags y Menciones}
\begin{itemize}
    \item Extraer hashtags de estilo Twitter/Instagram
    \item Extraer menciones (@usuario)
    \item Validar que no empiecen con números
\end{itemize}

\textbf{3. Limpiador de Texto HTML}
\begin{itemize}
    \item Remover todas las etiquetas HTML
    \item Preservar el contenido entre etiquetas
    \item Decodificar entidades HTML básicas
\end{itemize}

\subsection{Nivel Avanzado (+15 puntos cada uno)}

\textbf{4. Extractor de Información de Facturas}
\begin{itemize}
    \item Procesar texto de facturas costarricenses
    \item Extraer: número de factura, fecha, RUC, total
    \item Generar JSON estructurado
\end{itemize}

\textbf{5. Validador de Código de Programación}
\begin{itemize}
    \item Detectar y validar sintaxis Python básica
    \item Extraer nombres de funciones y variables
    \item Identificar imports
\end{itemize}

\textbf{6. Analizador de Sentimientos Básico}
\begin{itemize}
    \item Identificar palabras positivas/negativas usando regex
    \item Detectar intensificadores (muy, super, etc.)
    \item Generar score de sentimiento
\end{itemize}

\subsection{Nivel Experto (+20 puntos)}

\textbf{7. Parser de Lenguaje Natural Simple}
\begin{itemize}
    \item Extraer sujeto, verbo, objeto de oraciones simples
    \item Usar grupos nombrados extensivamente
    \item Manejar múltiples estructuras gramaticales
\end{itemize}

\textbf{8. Validador y Formateador de Direcciones CR}
\begin{itemize}
    \item Parsear direcciones costarricenses con regex
    \item Extraer provincia, cantón, distrito, señas
    \item Normalizar y formatear según estándar
\end{itemize}

\section{Recursos y Referencias}

\subsection{Documentación}

\begin{itemize}
    \item Módulo re de Python: \url{https://docs.python.org/3/library/re.html}
    \item Tutorial de regex: \url{https://docs.python.org/3/howto/regex.html}
    \item Regex101 (probador online): \url{https://regex101.com/}
    \item RegExr (herramienta visual): \url{https://regexr.com/}
\end{itemize}

\subsection{Consejos para el Desarrollo}

\begin{enumerate}
    \item \textbf{Prueba tus regex}: Usa herramientas online antes de implementar
    \item \textbf{Empieza simple}: Construye patrones incrementalmente
    \item \textbf{Usa raw strings}: Siempre prefiere \texttt{r''} para patrones
    \item \textbf{Documenta patrones}: Los regex pueden ser crípticos
    \item \textbf{Casos de prueba}: Ten ejemplos válidos e inválidos
    \item \textbf{Compila patrones}: Si los usas múltiples veces
\end{enumerate}

\section{Criterios de Entrega}

\subsection{Archivos Requeridos}

\begin{enumerate}
    \item \textbf{validador\_extractor.py}: Código principal
    \item \textbf{pruebas.txt}: Archivo con casos de prueba
    \item \textbf{logs\_ejemplo.txt}: Archivo de log para procesar
    \item \textbf{README.txt}: Documentación del proyecto
\end{enumerate}

\subsection{Contenido del README}

\begin{itemize}
    \item Nombre del estudiante
    \item Descripción de patrones regex implementados
    \item Instrucciones de uso
    \item Ejemplos de validación
    \item Desafíos opcionales completados
    \item Limitaciones conocidas
\end{itemize}

\vspace{1cm}

\begin{center}
\textit{"Las expresiones regulares son como el kung-fu:}\\
\textit{parecen mágicas hasta que comprendes los fundamentos."}\\
\vspace{0.3cm}
--- Adaptación de Jeff Atwood
\end{center}

\vspace{0.5cm}

\begin{center}
\textbf{¡Éxito con tus expresiones regulares!}\\
\textit{Recuerda: Un buen patrón regex puede ahorrarte miles de líneas de código.}
\end{center}

\end{document}
